% Path potential and supply (shared)

\begin{lemma}[Path potential and supply]\label{lem:path-lb}
Let $P=(v_0,v_1,\ldots,v_k)$ be a path and let $t=v_k$.
Define
\[
\Psi(p)=\sum_{i=0}^{k} p(v_i)\,2^{-(k-i)}.
\]
Then $\Psi(p)$ is non-increasing under every pebbling move on~$P$.
In particular, if $p(v_0)=2^k-1$ and $p(v_i)=0$ for $i>0$, then $t$ cannot be pebbled.

If $p(t)=0$ and some sequence from a configuration~$p$ pebbles~$t$, then some
reachable configuration~$p'$ satisfies $p'(v_0)\ge 2^k$.
\end{lemma}

\begin{proof}
A move $v_i\to v_{i+1}$ (toward~$t$) satisfies
$\Delta\Psi = -2\cdot 2^{-(k-i)}+2^{-(k-i-1)}=0$.
A move $v_{i+1}\to v_i$ satisfies $\Delta\Psi=-3\cdot 2^{-(k-i)}<0$.
Thus $\Psi$ never increases, proving the first part.

For the supply statement, some move in the pebbling sequence is the first to place
a pebble on~$t$; it must be $v_{k-1}\to t$ when $k\ge 1$.
Tracing backward along~$P$ from that first arrival, each pebble placed on~$v_j$
requires two pebbles on~$v_{j-1}$ at an earlier stage, so at least $2^k$ pebbles
must have appeared on~$v_0$ at some reachable configuration.
\end{proof}
